% SPDX-License-Identifier: GPL-3.0-or-later
% Copyright (C) 2026 Alessandro Greco
\chapter{Реальные атаки на сеть}\label{cap:27-casi-rete}
\aggiornatoal{23 сентября 2026 г.}{публикации Tor Project (2014, 2022, 2023); tor 0.4.9.13; лабораторная сеть}

\begin{obiettivi}
К концу этой главы вы сможете:
\begin{itemize}[nosep]
  \item восстановить атаку 2014 года с использованием ячеек \cella{RELAY\_EARLY}, объяснить добавленную в tor защиту и проверить её в работающей лаборатории (Лабораторная работа~17);
  \item описать инцидент KAX17 и последовавшую работу по наблюдению за ретрансляторами;
  \item связать DoS-атаки 2022--2023 годов с защитами, описанными в \cref{cap:15-dos};
  \item сформулировать общие выводы о защите открытой сети.
\end{itemize}
\end{obiettivi}

В предыдущей главе атаки рассматривались как исследовательские темы: модели, эксперименты и проценты. Здесь мы разберём три инцидента безопасности реальной сети по публичным сообщениям Tor Project. В случае KAX17 не установлены ни намерения оператора, ни факт конкретной вредоносной деятельности. Нас интересует прежде всего то, чему они учат о протоколе и об управлении открытой сетью, а не практические детали проведения атак.

% ---------------------------------------------------------------------------
\section{2014: сигнализация через типы ячеек}\label{sec:casi-relayearly}

\paragraph{Что произошло.} 4 июля 2014 года Tor Project обнаружил и исключил группу ретрансляторов, которые, по его мнению, пытались деанонимизировать людей, публиковавших дескрипторы onion-сервисов или запрашивавших сведения о них, тогда называвшиеся \emph{hidden services}. Ретрансляторы присоединились к сети 30 января: около 115 быстрых ретрансляторов без функции выхода в двух сетях /16, вместе составлявших примерно 6,4\,\% пропускной способности сторожевых узлов. 30 июля Tor Project опубликовал предупреждение о безопасности с описанием атаки и контрмер~\cite{torblog2014relayearly}.

\paragraph{Как работала атака.} Атака сочетала Sybil-размещение (\cref{sec:att-sybil}) с активным подтверждением трафика (\cref{sec:att-correlazione}). Ретрансляторы оставались стабильными достаточно долго, чтобы получить флаги \texttt{HSDir} и \texttt{Guard}. Когда вредоносный HSDir получал запрос дескриптора, он отвечал последовательностью ячеек, в которой чередование типов \cella{RELAY} и \cella{RELAY\_EARLY} кодировало имя запрошенного сервиса (\cref{fig:relayearly}). Если другой участник той же группы одновременно был сторожевым узлом этого клиента, он мог прочитать последовательность и связать клиента с запросом дескриптора сервиса, но не доказать, что он действительно посетил этот сервис.

\begin{figure}[htbp]
\centering
% SPDX-License-Identifier: GPL-3.0-or-later
% Copyright (C) 2026 Alessandro Greco
%% Русское альтернативное описание EPUB.
\altfig{Историческая атака 2014 года: вредоносный HSDir кодирует имя onion-сервиса типами ячеек RELAY и RELAY\_EARLY, направленных к клиенту. Сговорившийся с ним вредоносный guard читает видимые ретрансляторам типы и связывает клиента с запросом дескриптора, а не с доказанным посещением сервиса. Четыре нарисованные ячейки — лишь схема; на каналах заголовки защищены TLS.}
% Атака 2014 года: сигнал в типе ячейки, читаемый guard
\begin{tikzpicture}[font=\sffamily\scriptsize,
  ostile/.style={relay,fill=avviso!20,draw=avviso,align=center,font=\sffamily\scriptsize}]
  \node[client] (c) at (0,0) {клиент};
  \node[ostile] (g) at (3,0) {вредоносный\\guard};
  \node[relay] (m) at (6,0) {middle};
  \node[ostile] (h) at (9,0) {вредоносный\\HSDir};
  \draw[<-,thick] (c) -- (g);
  \draw[<-,thick] (g) -- (m);
  \draw[<-,thick] (m) -- (h);
  \foreach \x/\t/\c in {7.1/E/avviso!40,7.5/R/white,7.9/R/white,8.3/E/avviso!40} {
    \node[draw,minimum size=3.5mm,inner sep=0pt,fill=\c] at (\x,0.55) {\t};
  }
  \node[align=center,text=torgrigio] at (7.6,1.12) {схема: четыре ячейки для примера};
  \node[align=center,text=avviso] at (9,-1) {кодирует имя сервиса\\типами RELAY (R) и\\RELAY\_EARLY (E)};
  \node[align=center,text=avviso] at (3,-1) {читает типы ячеек:\\связывает клиента с\\запросом дескриптора};
  \node[align=center,text=torgrigio] at (6,-1) {передаёт типы ячеек\\как в атаке 2014 года};
\end{tikzpicture}

\caption{Историческая атака 2014 года: вредоносный HSDir кодировал запрошенный сервис внешними командами ячеек RELAY и RELAY\_EARLY, видимыми ретрансляторам в цепи, но защищёнными TLS на сетевых соединениях. Сговорившийся с ним guard мог связать клиента с запросом дескриптора, но не доказать посещение сервиса. Четыре ячейки показаны схематично и не кодируют полный onion-адрес.}
\label{fig:relayearly}
\end{figure}

Почему именно это поле? Ячейки \cella{RELAY\_EARLY} ограничивают длину цепи: клиенты обязаны использовать их для запросов расширения, а ретрансляторы принимают не более восьми на цепь (\spec{tor-spec}{Handling RELAY\_EARLY cells}). Однако старые проверки применялись только к трафику от клиента в сеть. Обычно никто не отправлял такие ячейки в обратном направлении, и никто их там не проверял. В отличие от содержимого, защищённого послойным шифрованием (\cref{cap:10-cifratura-relay}), внешняя команда ячейки видна каждому ретранслятору. Сигнал в этом поле проходит транзитом без изменений. Предупреждение также отмечало, что сигнал кодировал имя сервиса в читаемом виде, поэтому ретранслятор, получивший такие ячейки, мог распознать сигнал. Постороннему наблюдателю с одними лишь записями сетевых пакетов пришлось бы сначала расшифровать TLS между узлами, чтобы впоследствии извлечь его.

\paragraph{Ответ сети.} Последовали три шага: вредоносные ретрансляторы исключили; tor 0.2.4.23 и 0.2.5.6-alpha начали закрывать цепи при получении направленной к клиенту \cella{RELAY\_EARLY}; параллельно продолжалась работа по сокращению числа сторожевых узлов каждого клиента с трёх до одного, уменьшая число возможностей попасть в позицию наблюдения. Предупреждение подчёркивало, что \emph{пассивное} подтверждение трафика по-прежнему остаётся нерешённой проблемой. Канал через тип ячейки сделал подтверждение проще, а не впервые возможным.

% ---------------------------------------------------------------------------
\section{Лабораторная работа: лимит RELAY\_EARLY}\label{sec:lab-casi}

\begin{laboratorio}[17]{Восемь разрешены; девять — нет}
\textbf{Цель.} С помощью \file{labs/ru/код/torlab} проверить, как современный ретранслятор применяет лимит ячеек \cella{RELAY\_EARLY}.\\
\textbf{Предварительные условия.} Лабораторные работы~4 и~5; лабораторная сеть должна быть запущена.\\
\textbf{Оценочное время.} 30 минут.
\end{laboratorio}

Наш \file{circuit.py} уже отказывается отправлять более восьми ячеек \cella{RELAY\_EARLY}. Скрипт \file{labs/ru/лаборатория/relay\_early.py} намеренно обходит эту защиту. Он строит одноузловую цепь к \texttt{test004r}, отправляет $n$ ячеек \cella{RELAY\_EARLY}, содержащих сообщение \texttt{DROP}, которое ретранслятор отбрасывает без ответа, а затем запрашивает соединение с каталогом через \texttt{BEGIN\_DIR}, чтобы проверить, остаётся ли цепь активной:
\begin{lstlisting}[language=bash,numbers=none]
$ PYTHONPATH=labs/ru/код python3 labs/ru/лаборатория/relay_early.py 0 8 9
0 ячеек RELAY_EARLY -> CONNECTED
8 ячеек RELAY_EARLY -> CONNECTED
9 ячеек RELAY_EARLY -> circuit уничтожена, причина 1
\end{lstlisting}
Девятая ячейка закрывает цепь: relay отправляет \cella{DESTROY} с причиной 1, \texttt{PROTOCOL}. Листинг получен в эталонном запуске: адреса, идентификаторы цепей и записи журнала различаются между запусками. Эталонный журнал ретранслятора:
\begin{lstlisting}[numbers=none]
[warn] Received too many RELAY_EARLY cells on circ 4045911407 from
       127.0.0.1:58048.  Closing circuit.
\end{lstlisting}
Проверка находится в \file{src/core/or/command.c}, до расшифрования ячейки. Несколькими строками выше расположен контроль, добавленный в 2014 году: направленная к клиенту \cella{RELAY\_EARLY} вызывает предупреждение и закрытие цепи. Комментарий в исходном коде объясняет, что когда-то такие ячейки возникали из-за старой ошибки tor, но ретрансляторы с этой ошибкой давно исчезли. Сегодня такая ячейка считается \enquote{неожиданной} и обрабатывается как атака. Лабораторная работа также входит в \texttt{make verifica-lab}.

% ---------------------------------------------------------------------------
\section{2020--2021: KAX17 и здоровье сети}\label{sec:casi-kax17}

Второй инцидент показывает менее заметную часть защиты: сначала противника нужно обнаружить. Tor Project рассказал об этом случае в апреле 2022 года в статье о своей команде по состоянию сети~\cite{torblog2022badrelays}.

В декабре 2020 года графики ретрансляторов показали необычный паттерн. Каждый месяц, точно в полночь первого числа, в сеть входили десятки ретрансляторов, а затем постепенно исчезали. Это напоминало ошибку в механизме перехода в спящий режим (\opt{AccountingMax}, \cref{sec:relay-config}). В сентябре 2021 года начали появляться группы ретрансляторов, которые, судя по признакам, принадлежали одному оператору, но не объявляли семейство и не указывали контактные данные. Центры каталогов исключали их, но появлялись замены. В конце октября анонимное сообщение помогло определить группу, позднее получившую название KAX17; её удалили. После этого прекратились и ежемесячные пики, и волны появления ретрансляторов, из-за чего Tor Project счёл правдоподобным, что за обоими явлениями стоял один оператор. Его личность и цели остаются неизвестными.

Статья необычно откровенно описывает запоздалую реакцию: инструменты мониторинга годами оставались без должного внимания; доверие между добровольцами, сотрудниками и центрами каталогов было слабым; организация не была структурирована для такой работы. Она также объясняет процедуру управления. Специалисты по мониторингу предлагают исключить ретранслятор, но решение большинством принимают центры каталогов, а не сотрудники Tor Project. Это распределённое доверие из \cref{cap:12-directory}, применённое к эксплуатации сети.

\begin{approfondimento}[Ретрансляторы без объявления семейства]
Первым тревожным сигналом в случае KAX17 был не выявленная эксплуатация уязвимости, а нарушение социальной нормы: многочисленные ретрансляторы, очевидно управляемые одним оператором, без объявленной family (\cref{sec:relay-famiglie}). Более общий вывод состоит в том, что атакующему иногда не нужно ломать протокол, если он может обойти соглашения, на которых основан выбор маршрута.
\end{approfondimento}

% ---------------------------------------------------------------------------
\section{2022--2023: отказ в обслуживании}\label{sec:casi-dos}

Начиная с лета 2022 года сеть месяцами испытывала различные атаки отказа в обслуживании. В феврале 2023 года Tor Project сообщил, что пользователи иногда не могли загружать страницы или обращаться к onion-сервисам, что тактика и цели атак менялись со временем и что личности атакующих невозможно было определить с уверенностью~\cite{torblog2023dos}. Onion-сервисы были особенно уязвимы: они не могли блокировать IP-адреса атакующих, потому что запросы приходили по анонимным цепям.

Долгосрочным ответом стала система доказательства выполнения работы (proof of work), появившаяся в августе 2023 года в tor 0.4.8 (\cref{sec:dos-onion})~\cite{torblog2023pow}. В обычном состоянии защита неактивна; при атаке сервис просит клиентов решать вычислительные задачи возрастающей сложности и отдаёт приоритет тем, кто выполнил больше работы. Tor Project оценивал обычные затраты для нормального клиента в несколько миллисекунд, а при интенсивной атаке — потенциально примерно в минуту. Принцип тот же, что в Лабораторной работе~9: сделать дешёвую для атакующего атаку дорогой, не требуя от пользователя раскрывать личность.

% ---------------------------------------------------------------------------
\section{Выводы}\label{sec:casi-lezioni}

\begin{itemize}
  \item \textbf{Каждое видимое поле может стать каналом.} Атака 2014 года не взламывала криптографию. Она использовала поле, которое relay могли читать и пересылать. Для каждого незашифрованного бита протокол анонимности должен задаваться вопросами: кто может его записать и кто может прочитать? Это одна из мотиваций более новой схемы шифрования ретрансляционных сообщений (\cref{sec:cgo}).
  \item \textbf{Sybil-атаки могут быть дешёвыми; мониторинг — нет.} Добавление ретрансляторов стоит сравнительно мало по сравнению с поддержкой инструментов, компетенций и доверия, необходимых для выявления скоординированного вредоносного поведения.
  \item \textbf{Защиты нужно оценивать во времени.} Использование одного сторожевого узла и увеличение срока его использования и появление vanguards уменьшают число возможностей атакующего с течением времени (\cref{sec:att-correlazione}).
  \item \textbf{Прозрачность — часть защиты.} Во всех трёх случаях Tor Project публиковал и то, что было известно, и то, что оставалось неизвестным, позволяя пользователям, операторам и исследователям оценивать риск.
\end{itemize}

% ---------------------------------------------------------------------------
\section*{Итоги}
\begin{itemize}
  \item В 2014 году Sybil-группа, чьи ретрансляторы выполняли роли сторожевых узлов и HSDir, кодировала в типах ячеек имена onion-сервисов, дескрипторы которых запрашивали клиенты. Сейчас tor закрывает цепи при получении направленной к клиенту \cella{RELAY\_EARLY}, а переход от трёх сторожевых узлов к одному, а позднее к двум в 2018 году, сократил возможности выбора вредоносного сторожевого узла (\cref{sec:guard}).
  \item Лимит в восемь ячеек \cella{RELAY\_EARLY} можно проверить в лабораторной сети: девятая уничтожает цепь.
  \item KAX17 был выявлен после месяцев косвенных признаков, включая ежемесячные пики и необъявленные семейства ретрансляторов, а также после анонимного сообщения. Решение об исключении ретрансляторов принимают центры каталогов.
  \item DoS-атаки 2022--2023 годов способствовали внедрению доказательства выполнения работы для onion-сервисов.
\end{itemize}

% ---------------------------------------------------------------------------
\section*{Упражнения}

\begin{esercizio}[1]\label{es:27-campo}
Почему сигнал в зашифрованном содержимом ячейки, записанный HSDir, не мог быть прочитан сторожевым узлом, тогда как сигнал во внешнем типе ячейки мог?
\end{esercizio}

\begin{esercizio}[1]\label{es:27-uno}
Объясните, как сокращение числа сторожевых узлов с трёх до одного уменьшает риск атаки вроде случая 2014 года и какой недостаток это создаёт.
\end{esercizio}

\begin{esercizio}[2]\label{es:27-verso}
Измените Лабораторную работу~17 так, чтобы показать: ячейка \cella{RELAY\_EARLY} учитывается в лимите, даже если она не содержит запрос на расширение цепи. Затем исследуйте код tor и определите, что происходит, когда \cella{RELAY\_EARLY} приходит в направлении клиента.
\end{esercizio}

\begin{esercizio}[2]\label{es:27-segnali}
Используя данные CollecTor, какие признаки могли бы показать, что несколько ретрансляторов принадлежат одному оператору без объявленного семейства? Предложите три критерия и обсудите ложные срабатывания.
\end{esercizio}

\begin{esercizio}[3]\label{es:27-canali}
Перечислите поля ячеек и ретрансляционных сообщений, видимые промежуточному узлу открытым текстом (\cref{cap:08-celle}). Для каждого обсудите, может ли вредоносный ретранслятор использовать его как канал к другому ретранслятору в той же цепи, и укажите существующую или предлагаемую защиту.
\end{esercizio}

% ---------------------------------------------------------------------------
\section*{Дополнительная литература}
\begin{itemize}
  \item Tor Project, \titolo{Tor security advisory: \enquote{relay early} traffic confirmation attack}~\cite{torblog2014relayearly}.
  \item Tor Project, \titolo{Malicious relays and the health of the Tor network}~\cite{torblog2022badrelays}.
  \item Tor Project, \titolo{Introducing Proof-of-Work Defense for Onion Services}~\cite{torblog2023pow}.
\end{itemize}
